\documentclass[11pt]{article}
\usepackage[T1]{fontenc}
\usepackage[utf8]{inputenc}
\usepackage[spanish,es-nodecimaldot,es-noshorthands]{babel}
\usepackage{lmodern,microtype,amsmath,graphicx}
\usepackage[a4paper,margin=20mm,headheight=15pt]{geometry}
\usepackage{array,booktabs,tabularx}
\usepackage{tikz}
\usetikzlibrary{positioning,arrows.meta}
\usepackage{xcolor,fancyhdr}
\definecolor{tinta}{HTML}{193045}
\definecolor{teal}{HTML}{116A61}
\definecolor{gris}{HTML}{4A5C6A}
\definecolor{suave}{HTML}{EAF4F3}
\usepackage[colorlinks=true,allcolors=teal,bookmarksnumbered=true]{hyperref}
\pagestyle{fancy}
\fancyhf{}
\fancyhead[L]{\small\color{gris} ROBÓTICA AVANZADA · UNITREE GO2}
\fancyhead[R]{\small\color{gris} PROYECTO FINAL · 2026-2}
\fancyfoot[R]{\small\color{gris}\thepage}
\renewcommand{\headrulewidth}{.3pt}
\setlength{\parindent}{0pt}
\setlength{\parskip}{.5em}
\setlength{\fboxsep}{8pt}
\renewcommand{\arraystretch}{1.18}
\newcolumntype{L}{>{\raggedright\arraybackslash}X}
\newcommand{\recuadro}[1]{\par\noindent\colorbox{suave}{\parbox{\dimexpr\linewidth-2\fboxsep}{#1}}\par}
\newcommand{\repo}{https://github.com/daca2022/go2-rl-deploy}

\hypersetup{pdftitle={Guion de 15 minutos - Unitree Go2}}
\begin{document}
\color{tinta}
{\color{teal}\small\bfseries EXPOSICIÓN ORAL · PROYECTO FINAL\par}
\vspace{.65em}
{\fontsize{30}{33}\selectfont\bfseries Go2 en 15 minutos\par}
\vspace{.35em}
{\large Trote en MuJoCo, control PD y comparación de entrenamientos\par}

\section*{Secuencia de la exposición}
\begin{tabularx}{\linewidth}{@{}p{.13\linewidth}p{.25\linewidth}L@{}}
\toprule
\textbf{Tiempo} & \textbf{Rúbrica} & \textbf{Qué mostrar y decir}\\
\midrule
0:00–1:30 & Introducción · 1 punto & Plantea el objetivo. Coordinar 12 articulaciones y seguir cuatro consignas manteniendo el equilibrio. Pregunta si un currículo preparado para mayor dificultad mejora la política.\\
1:30–4:00 & Marco teórico · 2 puntos & Usa la Figura 2. El actor genera objetivos y el PD regula el movimiento de las juntas. Explica qué aprende PPO, qué recuerda la GRU y qué simula MuJoCo.\\
4:00–7:30 & Metodología · 3 puntos & Muestra el Draw.io. E1 y E5 usan 3.072.000 transiciones en plano. Explica las pruebas internas congeladas, las siete condiciones y el significado de aprobación, caída y RMSE.\\
7:30–13:00 & Resultados · 3 puntos & Reproduce los 20 s del video A0. Muestra las Figuras 3 y 4. Explica el margen de E1 con masa y retardo y la mejora de E5 en terreno rugoso. Si usas demo en vivo, elige el modo MuJoCo del panel A2.\\
13:00–15:00 & Conclusiones · 1 punto & Resume qué funciona, dónde fallan ambos y qué experimento separaría el efecto de la pendiente del efecto del entorno.\\
\bottomrule
\end{tabularx}

\section*{Guía para la demostración}
\textbf{Video listo para reproducir.} Abre \texttt{Demo-A0-U1600-20s.mp4} antes de la exposición. Durante la reproducción identifica los apoyos diagonales y la consigna de avance de 0,30 m/s. El clip es A0/P0 U1600 y los resultados comparativos del informe son de U2000.

\textbf{Opción en vivo.} En el repositorio \texttt{go2-rl-deploy} ya instalado, ejecuta \texttt{bash abrir-panel.sh}. Selecciona el modo MuJoCo y la consigna Avance. El panel ejecuta A2/P0 y guarda el CSV de la prueba. Ten preparado el video local si el entorno gráfico tarda en abrir.

\section*{Preguntas que conviene ensayar}
\textbf{¿Por qué E5 tomó más horas?} Dedicó unas cuatro horas a pruebas internas con pesos congelados. Los dos entrenaron con el mismo número de transiciones.

\textbf{¿Por qué E1 gana en masa y retardo?} En avance, su RMSE nominal es 0,0608 m/s y sube a 0,0917. El de E5 va de 0,0743 a 0,1032. Los aumentos son parecidos y solo E5 supera el límite de 0,10.

\textbf{¿Más actualizaciones aseguraron mejor A2?} U2800 bajó ligeramente el error lateral en la revisión interna. U2000 conservó el criterio completo de trote al retroceder en pendiente leve.

\href{\repo}{Repositorio público y panel}\quad
\href{https://drive.google.com/drive/folders/1S8uEebuWrVEXju7h-1b8J_P_9BeeqBDS}{Material de la exposición}
\end{document}
